\documentclass[11pt,a4paper]{article}
\usepackage{fontspec,amsmath,amssymb,amsthm,microtype}
\setmainfont{lmroman10-regular.otf}[BoldFont=lmroman10-bold.otf,ItalicFont=lmroman10-italic.otf,BoldItalicFont=lmroman10-bolditalic.otf]
\usepackage{xcolor}
\usepackage[most]{tcolorbox}
\usepackage{mathrsfs}
\providecommand{\lt}{<}
\providecommand{\gt}{>}
\setmonofont{DejaVuSansMono.ttf}[BoldFont=DejaVuSansMono-Bold.ttf,ItalicFont=DejaVuSansMono-Oblique.ttf,BoldItalicFont=DejaVuSansMono-BoldOblique.ttf,Scale=0.86]
\newfontfamily\leanfallback{STIXTwoMath-Regular.otf}[Scale=0.86]
\definecolor{leanrule}{gray}{0.62}
\definecolor{leanback}{gray}{0.962}
\definecolor{leantitle}{gray}{0.22}
\definecolor{leanmuted}{gray}{0.38}
\newtcolorbox{leanbox}[1]{enhanced,breakable,colback=leanback,colframe=leanrule,
  boxrule=0.4pt,arc=0.6mm,left=2.4mm,right=2.4mm,top=1.8mm,bottom=1.8mm,
  before skip=0.8em,after skip=1em,fonttitle=\small\bfseries,coltitle=leantitle,
  colbacktitle=leanback,titlerule=0.3pt,title={#1}}
\newlength{\leancharwidth}
\newenvironment{leancode}{\par\smallskip\begingroup\ttfamily\footnotesize\raggedright
  \settowidth{\leancharwidth}{M}\setlength{\parskip}{0pt}\setlength{\parindent}{0pt}
  \frenchspacing}{\par\endgroup\smallskip}
\newcommand{\leanline}[2]{\hangindent=\dimexpr#1\leancharwidth+4\leancharwidth\relax
  \hangafter=1\noindent\hspace*{#1\leancharwidth}#2\par}
\newcommand{\leansource}[1]{\par{\footnotesize\color{leanmuted}#1}\par}
\newcommand{\leanlabel}[1]{\par\smallskip\noindent{\small\bfseries #1}\par}
\usepackage[margin=25mm]{geometry}
\usepackage[unicode,hidelinks]{hyperref}
\hypersetup{pdfauthor={LangevinLean Project}}
\setlength{\emergencystretch}{3em}
\setlength{\parskip}{0.4em}
\setlength{\parindent}{0pt}
\allowdisplaybreaks[2]
\raggedbottom
\title{\texorpdfstring{Exact algebraic corrections and an averaged-proximity counterexample}{Exact algebraic corrections and an averaged-proximity counterexample}}
\author{LangevinLean Project}
\date{}
\begin{document}\maketitle
Draft mathematical exposition; these notes may contain mistakes.


The examples below establish specific algebraic facts used to qualify the source arguments: the sign of a quadratic proximal expansion, scaling of a squared norm, the inverse-step factor in residual monotonicity, and failure of shared-state proximal optimality for averaged pairwise outputs. The convergence results use the finite potential-gap and summed energy inequalities proved above. Each example has only the scope stated for its own assumptions.



\section{\texorpdfstring{The second-order sign in the quadratic proximal expansion}{The second-order sign in the quadratic proximal expansion}}


The proximal equation \(p+\lambda\nabla g_2(p)=x\) gives a positive second-order coefficient in the scalar quadratic example:

\[
\operatorname{prox}_{u^2/2}^\lambda(x)
=\frac{x}{1+\lambda}
=x-\lambda x+\lambda^2x-\frac{\lambda^3x}{1+\lambda}.
\]

The negative coefficient printed in the source leaves error at least \(\lambda^2\) at \(x=1\) and therefore cannot be its third-order expansion. The quadratic is not globally value Lipschitz, so this calculation diagnoses the displayed algebra rather than refuting a full theorem under that assumption.



Almost-everywhere twice differentiability and the stated moment condition also do not provide the uniform integrable third-order remainder needed for the claimed order \(O(\lambda)\) for displacement of the original marginal maximizer. The proved finite-gap result gives \(\sqrt{\lambda K^2/m}\); this difference in established scope does not prove the stronger assertion false. For \(f(u)=u^2/2\) the minimizer equation is \((1+\lambda)p=x\). Division gives the exact resolvent and its positive \(\lambda^2x\) term with the displayed third-order remainder.

\begin{leanbox}{Lean: \texttt{quadratic\_\allowbreak{}prox\_\allowbreak{}minimizes}, \texttt{quadratic\_\allowbreak{}prox\_\allowbreak{}expansion}, \texttt{negative\_\allowbreak{}second\_\allowbreak{}order\_\allowbreak{}residual} and 1 more}
\leansource{\texttt{quadratic\_\allowbreak{}prox\_\allowbreak{}minimizes} \quad ProximalLangevin/\allowbreak{}Corrections.lean, lines 20–31}
\begin{leancode}
\leanline{0}{theorem\ quadratic\_prox\_minimizes\ (x\ lam\ :\ ℝ)\ (hlam\ :\ 0\ <\ lam)\ (u\ :\ ℝ)\ :}
\leanline{4}{(x\ /\ (1\ +\ lam))\ \textasciicircum{}\ 2\ /\ 2\ +\ (x\ /\ (1\ +\ lam)\ -\ x)\ \textasciicircum{}\ 2\ /\ (2\ *\ lam)\ ≤}
\leanline{6}{u\ \textasciicircum{}\ 2\ /\ 2\ +\ (u\ -\ x)\ \textasciicircum{}\ 2\ /\ (2\ *\ lam)}
\end{leancode}
\leansource{\texttt{quadratic\_\allowbreak{}prox\_\allowbreak{}expansion} \quad ProximalLangevin/\allowbreak{}Corrections.lean, lines 36–39}
\begin{leancode}
\leanline{0}{theorem\ quadratic\_prox\_expansion\ (x\ lam\ :\ ℝ)\ (hlam\ :\ 1\ +\ lam\ ≠\ 0)\ :}
\leanline{4}{x\ /\ (1\ +\ lam)\ =\ x\ -\ lam\ *\ x\ +\ lam\ \textasciicircum{}\ 2\ *\ x\ -\ lam\ \textasciicircum{}\ 3\ *\ x\ /\ (1\ +\ lam)}
\end{leancode}
\leansource{\texttt{negative\_\allowbreak{}second\_\allowbreak{}order\_\allowbreak{}residual} \quad ProximalLangevin/\allowbreak{}Corrections.lean, lines 43–46}
\begin{leancode}
\leanline{0}{theorem\ negative\_second\_order\_residual\ (lam\ :\ ℝ)\ (hlam\ :\ 0\ <\ lam)\ :}
\leanline{4}{1\ /\ (1\ +\ lam)\ -\ (1\ -\ lam\ -\ lam\ \textasciicircum{}\ 2)\ =\ lam\ \textasciicircum{}\ 2\ *\ (2\ +\ lam)\ /\ (1\ +\ lam)}
\end{leancode}
\leansource{\texttt{negative\_\allowbreak{}second\_\allowbreak{}order\_\allowbreak{}residual\_\allowbreak{}ge\_\allowbreak{}sq} \quad ProximalLangevin/\allowbreak{}Corrections.lean, lines 48–52}
\begin{leancode}
\leanline{0}{theorem\ negative\_second\_order\_residual\_ge\_sq\ (lam\ :\ ℝ)\ (hlam\ :\ 0\ <\ lam)\ :}
\leanline{4}{lam\ \textasciicircum{}\ 2\ ≤\ 1\ /\ (1\ +\ lam)\ -\ (1\ -\ lam\ -\ lam\ \textasciicircum{}\ 2)}
\end{leancode}
\leanlabel{How the Lean proof follows the argument}
Formal correspondence: quadratic\_prox\_minimizes : The proximal point of \(u\mapsto u^2/2\) is \(x/(1+\lambda)\). The difference of objectives is a nonnegative square; no differentiability of an argmin map is assumed. quadratic\_prox\_expansion : The second-order coefficient in the proximal expansion has a positive sign. This exact identity diagnoses the opposite sign in v1 equation (25) and v3 equation (22). negative\_second\_order\_residual : The negative second-order term printed in v1 equation (25) /\allowbreak{} v3 equation (22) leaves an error of order \(\lambda^2\), already for the scalar quadratic.



\texttt{quadratic\_\allowbreak{}prox\_\allowbreak{}minimizes}: the intermediate facts include \texttt{hd}, \texttt{hidentity}, \texttt{hnonneg}, in source order. The remaining scalar rearrangements are checked by \texttt{field\_\allowbreak{}simp}, \texttt{ring}, \texttt{linarith}.



\texttt{quadratic\_\allowbreak{}prox\_\allowbreak{}expansion}: the proof is the application \texttt{field\_\allowbreak{}simp}. It specializes that result with the hypotheses and objects in the exact statement.



\texttt{negative\_\allowbreak{}second\_\allowbreak{}order\_\allowbreak{}residual}: the proof is the application \texttt{field\_\allowbreak{}simp}. It specializes that result with the hypotheses and objects in the exact statement.



\texttt{negative\_\allowbreak{}second\_\allowbreak{}order\_\allowbreak{}residual\_\allowbreak{}ge\_\allowbreak{}sq}: the source steps apply or rewrite \texttt{negative\_\allowbreak{}second\_\allowbreak{}order\_\allowbreak{}residual}, \texttt{le\_\allowbreak{}div\_\allowbreak{}iff\ensuremath{{}_{0}}}. The remaining scalar rearrangements are checked by \texttt{nlinarith}.
\end{leanbox}

\section{\texorpdfstring{Squared scaling and the unit recurrence ratio}{Squared scaling and the unit recurrence ratio}}


Multiplication of a vector by \(N\) multiplies its squared norm by \(N^2\), by norm homogeneity and squaring. Separately, an affine recurrence with unit ratio has finite sum \(n\) and therefore a linear-in-iteration additive bound, rather than a bounded geometric sum. These identities determine the particle powers and the strict positive-denominator condition needed in an energy or transport estimate.

\begin{leanbox}{Lean: \texttt{squared\_\allowbreak{}gradient\_\allowbreak{}scaling}, \texttt{linear\_\allowbreak{}squared\_\allowbreak{}gradient\_\allowbreak{}scaling\_\allowbreak{}fails}, \texttt{endpoint\_\allowbreak{}denominator\_\allowbreak{}zero} and 1 more}
\leansource{\texttt{squared\_\allowbreak{}gradient\_\allowbreak{}scaling} \quad ProximalLangevin/\allowbreak{}Corrections.lean, lines 56–59}
\begin{leancode}
\leanline{0}{theorem\ squared\_gradient\_scaling\ \{E\ :\ Type*\}\ [NormedAddCommGroup\ E]}
\leanline{4}{[NormedSpace\ ℝ\ E]\ (N\ :\ ℝ)\ (g\ :\ E)\ :}
\leanline{4}{‖N\ •\ g‖\ \textasciicircum{}\ 2\ =\ N\ \textasciicircum{}\ 2\ *\ ‖g‖\ \textasciicircum{}\ 2}
\end{leancode}
\leansource{\texttt{linear\_\allowbreak{}squared\_\allowbreak{}gradient\_\allowbreak{}scaling\_\allowbreak{}fails} \quad ProximalLangevin/\allowbreak{}Corrections.lean, lines 63–64}
\begin{leancode}
\leanline{0}{theorem\ linear\_squared\_gradient\_scaling\_fails\ :}
\leanline{4}{¬\ ((2\ *\ (1\ :\ ℝ))\ \textasciicircum{}\ 2\ ≤\ 2\ *\ 1)}
\end{leancode}
\leansource{\texttt{endpoint\_\allowbreak{}denominator\_\allowbreak{}zero} \quad ProximalLangevin/\allowbreak{}Corrections.lean, lines 69–73}
\begin{leancode}
\leanline{0}{theorem\ endpoint\_denominator\_zero\ (μ\ γ\ :\ ℝ)}
\leanline{4}{(hγ\ :\ γ\ ≠\ 0)\ (hden\ :\ 1\ -\ μ\ *\ γ\ ≠\ 0)\ :}
\leanline{4}{1\ -\ (γ\ /\ (1\ -\ μ\ *\ γ))\ *\ (1\ -\ μ\ *\ γ)\ /\ γ\ =\ 0}
\end{leancode}
\leansource{\texttt{endpoint\_\allowbreak{}error\_\allowbreak{}floor\_\allowbreak{}fails} \quad ProximalLangevin/\allowbreak{}Corrections.lean, lines 77–78}
\begin{leancode}
\leanline{0}{theorem\ endpoint\_error\_floor\_fails\ :}
\leanline{4}{¬\ ((1\ :\ ℝ)\ ≤\ 0\ +\ 1\ /\ (1\ -\ 1))}
\end{leancode}
\leanlabel{How the Lean proof follows the argument}
Norm homogeneity gives \(\|Ng\|^2=N^2\|g\|^2\). The numerical case \(N=2,\|g\|=1\) disproves linear squared-gradient scaling. Substituting \(\lambda=\gamma/(1-\mu\gamma)\) gives \(1-q=0\). Lean’s total real division has \(1/0=0\); the numerical certificate \(1\not\le0+1/(1-1)\) shows why an unguarded formula at that endpoint cannot express an error floor.



Formal correspondence: squared\_gradient\_scaling : Scaling a gradient by \(N\) scales its squared norm by \(N^2\), not \(N\). This corrects v1 Corollary A.5; v3 Corollary A.12 already uses this scaling. linear\_squared\_gradient\_scaling\_fails : A numerical witness to the invalid \(NC\) bound from the valid \(C\) bound when the gradient is multiplied by \(N\). endpoint\_denominator\_zero : The geometric-series denominator is zero at the upper endpoint allowed in v1 Corollaries A.4--A.6 and v3 Corollaries A.11--A.13. A finite error floor requires a strict inequality. endpoint\_error\_floor\_fails : At \(q=1\), the recurrence \(e(n+1)=e(n)+1\) has linear growth. Interpreting real division by zero cannot rescue the printed error floor.



\texttt{squared\_\allowbreak{}gradient\_\allowbreak{}scaling}: the source steps apply or rewrite \texttt{norm\_\allowbreak{}smul}, \texttt{Real.\allowbreak{}norm\_\allowbreak{}eq\_\allowbreak{}abs}, \texttt{mul\_\allowbreak{}pow}, \texttt{sq\_\allowbreak{}abs}.



\texttt{linear\_\allowbreak{}squared\_\allowbreak{}gradient\_\allowbreak{}scaling\_\allowbreak{}fails}: the proof is the application \texttt{norm\_\allowbreak{}num}. It specializes that result with the hypotheses and objects in the exact statement.



\texttt{endpoint\_\allowbreak{}denominator\_\allowbreak{}zero}: the source steps apply or rewrite \texttt{div\_\allowbreak{}mul\_\allowbreak{}cancel\ensuremath{{}_{0}}}, \texttt{div\_\allowbreak{}self}.



\texttt{endpoint\_\allowbreak{}error\_\allowbreak{}floor\_\allowbreak{}fails}: the proof is the application \texttt{norm\_\allowbreak{}num}. It specializes that result with the hypotheses and objects in the exact statement.
\end{leanbox}

\section{\texorpdfstring{The inverse-step factor in residual monotonicity}{The inverse-step factor in residual monotonicity}}


Assume \(\langle x-y,p(x)-p(y)\rangle\le(1-\mu)\|x-y\|^2\) and define the residual difference by \(R(x)-R(y)=\gamma^{-1}[(x-y)-(p(x)-p(y))]\). For \(\gamma>0\), expansion gives \(\langle x-y,R(x)-R(y)\rangle\ge(\mu/\gamma)\|x-y\|^2\). Hence the strong-monotonicity coefficient is \(\mu/\gamma\) rather than \(\mu\).



Explicitly the inner product expands as \(\langle x-y,R(x)-R(y)\rangle=\gamma^{-1}(\|x-y\|^2-\langle x-y,p(x)-p(y)\rangle)\). Substitution of the proximal deficit bounds it below by \((\mu/\gamma)\|x-y\|^2\). Division by \(\gamma\) in the residual definition is the source of this inverse-step factor.

\begin{leanbox}{Lean: \texttt{proximal\_\allowbreak{}displacement\_\allowbreak{}strong\_\allowbreak{}monotonicity}}
\leansource{\texttt{proximal\_\allowbreak{}displacement\_\allowbreak{}strong\_\allowbreak{}monotonicity} \quad ProximalLangevin/\allowbreak{}Corrections.lean, lines 82–94}
\begin{leancode}
\leanline{0}{theorem\ proximal\_displacement\_strong\_monotonicity}
\leanline{4}{\{E\ :\ Type*\}\ [NormedAddCommGroup\ E]\ [InnerProductSpace\ ℝ\ E]}
\leanline{4}{(p\ :\ E\ →\ E)\ (μ\ γ\ :\ ℝ)\ (hγ\ :\ 0\ <\ γ)}
\leanline{4}{(hB\ :\ ∀\ x\ y,\ inner\ ℝ\ (x\ -\ y)\ (p\ x\ -\ p\ y)\ ≤\ (1\ -\ μ)\ *\ ‖x\ -\ y‖\ \textasciicircum{}\ 2)}
\leanline{4}{(x\ y\ :\ E)\ :}
\leanline{4}{(μ\ /\ γ)\ *\ ‖x\ -\ y‖\ \textasciicircum{}\ 2\ ≤}
\leanline{6}{inner\ ℝ\ (x\ -\ y)\ (γ⁻¹\ •\ (x\ -\ p\ x)\ -\ γ⁻¹\ •\ (y\ -\ p\ y))}
\end{leancode}
\leanlabel{How the Lean proof follows the argument}
Formal correspondence: proximal\_displacement\_strong\_monotonicity : If the proximal displacement satisfies the paper's assumption B1, the envelope gradient inherits the constant \(\mu/\gamma\).



\texttt{proximal\_\allowbreak{}displacement\_\allowbreak{}strong\_\allowbreak{}monotonicity}: the intermediate facts include \texttt{hid}, \texttt{h}, in source order; the source steps apply or rewrite \texttt{inner\_\allowbreak{}smul\_\allowbreak{}right}, \texttt{inner\_\allowbreak{}sub\_\allowbreak{}right}, \texttt{real\_\allowbreak{}inner\_\allowbreak{}self\_\allowbreak{}eq\_\allowbreak{}norm\_\allowbreak{}sq}, \texttt{div\_\allowbreak{}eq\_\allowbreak{}mul\_\allowbreak{}inv}. The remaining scalar rearrangements are checked by \texttt{nlinarith}.
\end{leanbox}

\section{\texorpdfstring{A counterexample to shared-state proximal optimality}{A counterexample to shared-state proximal optimality}}


Take \(g(\theta,x)=|\theta-x|\) and proximal scale one. The pair inputs \((1,0)\) and \((1,4)\) have proximal minimizers \((1/2,1/2)\) and \((2,3)\). Averaging their parameter components gives \((5/4,1/2,3)\). The shared-state objective is

\[\tfrac12(|\theta-x_1|+|\theta-x_2|)+\tfrac12(\theta-1)^2+\tfrac14(x_1^2+(x_2-4)^2).\]

The competitor \((1,1,3)\) has strictly smaller objective, as follows from direct absolute-value evaluation. Each pair was allowed its own proposed parameter, whereas the shared objective permits only one. Averaging the proposals need not solve that common minimization problem.



In particular the averaged output has value \(51/32\) and the feasible comparator has value \(3/2=48/32\). This strict comparison disproves the shared-state proximal identity in general. The splitting proof instead sums the valid single-pair energy inequalities and applies Jensen’s inequality to the parameter average.

\begin{leanbox}{Lean: \texttt{pairAbsObjective}, \texttt{pairAbs\_\allowbreak{}first\_\allowbreak{}prox}, \texttt{pairAbs\_\allowbreak{}second\_\allowbreak{}prox} and 2 more}
\leansource{\texttt{pairAbsObjective} \quad ProximalLangevin/\allowbreak{}Corrections.lean, lines 98–99}
\begin{leancode}
\leanline{0}{def\ pairAbsObjective\ (θ\ x\ θ₀\ x₀\ :\ ℝ)\ :\ ℝ\ :=}
\leanline{2}{|θ\ -\ x|\ +\ ((θ\ -\ θ₀)\ \textasciicircum{}\ 2\ +\ (x\ -\ x₀)\ \textasciicircum{}\ 2)\ /\ 2}
\end{leancode}
\leansource{\texttt{pairAbs\_\allowbreak{}first\_\allowbreak{}prox} \quad ProximalLangevin/\allowbreak{}Corrections.lean, lines 102–107}
\begin{leancode}
\leanline{0}{theorem\ pairAbs\_first\_prox\ (θ\ x\ :\ ℝ)\ :}
\leanline{4}{pairAbsObjective\ (1\ /\ 2)\ (1\ /\ 2)\ 1\ 0\ ≤\ pairAbsObjective\ θ\ x\ 1\ 0}
\end{leancode}
\leansource{\texttt{pairAbs\_\allowbreak{}second\_\allowbreak{}prox} \quad ProximalLangevin/\allowbreak{}Corrections.lean, lines 110–114}
\begin{leancode}
\leanline{0}{theorem\ pairAbs\_second\_prox\ (θ\ x\ :\ ℝ)\ :}
\leanline{4}{pairAbsObjective\ 2\ 3\ 1\ 4\ ≤\ pairAbsObjective\ θ\ x\ 1\ 4}
\end{leancode}
\leansource{\texttt{sharedAbsObjective} \quad ProximalLangevin/\allowbreak{}Corrections.lean, lines 118–120}
\begin{leancode}
\leanline{0}{def\ sharedAbsObjective\ (θ\ x₁\ x₂\ :\ ℝ)\ :\ ℝ\ :=}
\leanline{2}{(|θ\ -\ x₁|\ +\ |θ\ -\ x₂|)\ /\ 2\ +}
\leanline{4}{((θ\ -\ 1)\ \textasciicircum{}\ 2\ +\ (x₁\ \textasciicircum{}\ 2\ +\ (x₂\ -\ 4)\ \textasciicircum{}\ 2)\ /\ 2)\ /\ 2}
\end{leancode}
\leansource{\texttt{averaged\_\allowbreak{}pair\_\allowbreak{}prox\_\allowbreak{}not\_\allowbreak{}lifted\_\allowbreak{}minimizer} \quad ProximalLangevin/\allowbreak{}Corrections.lean, lines 126–128}
\begin{leancode}
\leanline{0}{theorem\ averaged\_pair\_prox\_not\_lifted\_minimizer\ :}
\leanline{4}{sharedAbsObjective\ 1\ 1\ 3\ <\ sharedAbsObjective\ (5\ /\ 4)\ (1\ /\ 2)\ 3}
\end{leancode}
\leanlabel{How the Lean proof follows the argument}
Formal correspondence: pairAbsObjective : The pair objective for the convex, globally Lipschitz penalty \(g_2(\theta,x)=|\theta-x|\) with smoothing parameter one. pairAbs\_first\_prox : At pair input \((1,0)\), the true proximal point is \((1/2,1/2)\). pairAbs\_second\_prox : At pair input \((1,4)\), the true proximal point is \((2,3)\). sharedAbsObjective : The two-particle lifted objective in physical particle coordinates. The quadratic penalty uses precisely the scaled state norm with \(N=2\). averaged\_pair\_prox\_not\_lifted\_minimizer : Averaging the two true pairwise proximal outputs gives parameter \(5/4\) and particles \((1/2,3)\). A feasible competitor \((1,1,3)\) has strictly smaller lifted objective. Thus the pairwise average cannot be a lifted proximal minimizer, even for a convex globally Lipschitz penalty.



\texttt{pair\allowbreak{}Abs\allowbreak{}Objective} is a definition with defining expression \texttt{|\ensuremath{\theta} -\allowbreak{} x| + ((\ensuremath{\theta} -\allowbreak{} \ensuremath{\theta}\ensuremath{{}_{0}}) \textasciicircum{} 2 + (x -\allowbreak{} x\ensuremath{{}_{0}}) \textasciicircum{} 2) /\allowbreak{} 2}. The right-hand side specifies the object used in the mathematical statement.



\texttt{pair\allowbreak{}Abs\_\allowbreak{}first\_\allowbreak{}prox}: the intermediate fact is \texttt{ha}, in source order. The remaining scalar rearrangements are checked by \texttt{nlinarith}.



\texttt{pair\allowbreak{}Abs\_\allowbreak{}second\_\allowbreak{}prox}: the intermediate fact is \texttt{ha}, in source order. The remaining scalar rearrangements are checked by \texttt{nlinarith}.



\texttt{shared\allowbreak{}Abs\allowbreak{}Objective} is a definition with defining expression \texttt{(|\ensuremath{\theta} -\allowbreak{} x\ensuremath{{}_{1}}| + |\ensuremath{\theta} -\allowbreak{} x\ensuremath{{}_{2}}|) /\allowbreak{} 2 +}. The right-hand side specifies the object used in the mathematical statement.



\texttt{averaged\_\allowbreak{}pair\_\allowbreak{}prox\_\allowbreak{}not\_\allowbreak{}lifted\_\allowbreak{}minimizer}: the proof is the application \texttt{norm\_\allowbreak{}num [shared\allowbreak{}Abs\allowbreak{}Objective]}. It specializes that result with the hypotheses and objects in the exact statement.
\end{leanbox}

\end{document}
